\documentclass[11pt, letterpaper]{article}
\usepackage{amsmath, amssymb}
\usepackage{hyperref}
\usepackage{booktabs}
\usepackage{enumitem}
\usepackage{xcolor}
\usepackage{geometry}
\usepackage{fancyhdr}
\usepackage{titlesec}
\usepackage{tcolorbox}
\usepackage{parskip}
\geometry{margin=1in}

\definecolor{courseblue}{HTML}{002452}
\definecolor{coursegold}{HTML}{FABD0F}
\definecolor{coursered}{HTML}{B90E31}
\definecolor{lightblue}{HTML}{E8EEF7}

\hypersetup{colorlinks=true,linkcolor=courseblue,urlcolor=coursered}

\tcbuselibrary{skins}
\newtcolorbox{guidancebox}{enhanced,colback=lightblue,colframe=courseblue,
  sharp corners,leftrule=4pt,left=6pt,right=6pt,top=4pt,bottom=4pt}

\pagestyle{fancy}\fancyhf{}
\fancyhead[L]{\small\textcolor{courseblue}{\textbf{CISC 867 --- Deep Learning | Fall 2026}}}
\fancyhead[R]{\small\textcolor{courseblue}{Project Proposal}}
\fancyfoot[C]{\small\thepage}
\renewcommand{\headrulewidth}{0.4pt}

\titleformat{\section}{\large\bfseries\color{courseblue}}{}{0em}{}[\titlerule]
\titleformat{\subsection}{\normalsize\bfseries\color{courseblue}}{}{0em}{}

%% ════════════════════════════════════════════════════════════════
%% INSTRUCTIONS FOR STUDENTS
%% Replace every [PLACEHOLDER] with your actual content.
%% Delete all \guidancebox blocks before submitting.
%% Submit as a compiled PDF + .tex source in a .zip archive.
%% ════════════════════════════════════════════════════════════════

\begin{document}

{\centering
  {\Large\textbf{CISC 867: Deep Learning --- Project Proposal}}\\[0.5em]
  {\large\textit{[PROJECT ID AND TITLE, e.g., P08: Abstractive Summarization via LoRA]}}\\[0.4em]
  \textbf{Group Members:}\\{}
  [Student 1 Name] \quad [Student ID] \quad (Role: Project Lead)\\{}
  [Student 2 Name] \quad [Student ID] \quad (Role: Research Lead)\\{}
  [Student 3 Name] \quad [Student ID] \quad (Role: Engineering Lead)\\[0.3em]
  \textbf{GitHub Repository:} \url{[https://github.com/your-org/your-repo]}\\[0.3em]
  \textbf{Date:} \today\\[0.4em]
  \rule{\linewidth}{0.4pt}
}

\section{Problem Statement and Motivation}

\begin{guidancebox}
\textbf{Guidance (delete before submitting).}
Clearly define the problem you are solving. Explain \emph{why} it matters
in a broader context. Reference 1--2 high-level papers or application areas.
This section should answer: ``What problem are we addressing, and who cares?''
Target: $\approx$300 words.
\end{guidancebox}

[WRITE YOUR PROBLEM STATEMENT HERE. Describe the task, why it is hard, and why
it is important. Situate it within a research or application context.]

\section{Literature Survey}

\begin{guidancebox}
\textbf{Guidance (delete before submitting).}
Summarise at least \textbf{5 related papers}. For each paper:
(1) state the core idea in one sentence, (2) describe what it achieves,
and (3) explain how your project relates to or differs from it.
Avoid copying abstracts verbatim. Identify at least one open problem
that motivates your proposed approach.
\end{guidancebox}

\subsection{Related Works}

\paragraph{[Paper 1: Authors, Year, Venue.]}
[One-paragraph summary: what they proposed, what benchmark they evaluated on,
what metric they achieved, and how this relates to your project.]

\paragraph{[Paper 2: Authors, Year, Venue.]}
[Summary\ldots]

\paragraph{[Paper 3: Authors, Year, Venue.]}
[Summary\ldots]

\paragraph{[Paper 4: Authors, Year, Venue.]}
[Summary\ldots]

\paragraph{[Paper 5: Authors, Year, Venue.]}
[Summary\ldots]

\subsection{Open Problem}

\begin{guidancebox}
\textbf{Guidance.} Identify a specific gap, limitation, or unexplored direction
in the related work that motivates your novel contribution. Be precise.
\end{guidancebox}

[DESCRIBE THE OPEN PROBLEM OR GAP YOUR PROJECT WILL ADDRESS.]

\section{Proposed Methodology}

\begin{guidancebox}
\textbf{Guidance.}
Describe your approach in enough detail that it could be reproduced:
(1) baseline model and dataset you will start with,
(2) your proposed modification or extension and \emph{why} you chose it,
(3) how you will evaluate success (metrics, baselines).
Include a diagram if helpful. Target: $\approx$400 words.
\end{guidancebox}

\subsection{Baseline}

[DESCRIBE THE BASELINE MODEL AND DATASET. Which existing code/checkpoint will
you start from? Which metric will you use to evaluate the baseline?]

\subsection{Proposed Extension}

[DESCRIBE YOUR PROPOSED MODIFICATION OR EXTENSION. Why do you expect it to
help? What is the key hypothesis you are testing?]

\subsection{Evaluation Plan}

[DESCRIBE YOUR EVALUATION PLAN:
\begin{itemize}[noitemsep]
  \item Dataset(s) and split(s) you will use
  \item Metrics (e.g., PSNR/SSIM, ROUGE-L, F1, AUROC)
  \item Baselines you will compare against
  \item Number of seeds / runs for statistical reliability
\end{itemize}]

\subsection{Early Novel Contribution (Optional Bonus)}

\begin{guidancebox}
\textbf{Guidance.} If you have already identified a specific novel contribution
(a gap in existing evaluations, an untested configuration, a new domain),
describe it here for the +2 bonus points. Be concrete and specific.
\end{guidancebox}

[DESCRIBE YOUR NOVEL CONTRIBUTION IDEA, OR LEAVE THIS SECTION EMPTY.]

\section{Project Timeline and Roles}

\begin{guidancebox}
\textbf{Guidance.} Provide a week-by-week plan. Assign each task to a named
group member. The plan must be specific enough to hold yourselves accountable.
\end{guidancebox}

\begin{center}
\begin{tabular}{lp{7cm}l}
\toprule
\textbf{Week} & \textbf{Tasks} & \textbf{Responsible} \\
\midrule
1 & [Task, e.g., Finish reading all 4 core papers; repo setup] & [Name] \\
2 & [Task] \textbf{(Proposal due)} & [Name] \\
3 & [Task] & [Name] \\
4 & [Task] & [Name] \\
5 & [Task] & [Name] \\
6 & [Task] \textbf{(Midterm due)} & [Name] \\
7 & [Task] & [Name] \\
8 & [Task] & [Name] \\
9 & [Task] & [Name] \\
10 & [Task] \textbf{(Final due)} & [Name] \\
\bottomrule
\end{tabular}
\end{center}

\subsection{Contribution Declaration}

\begin{tabular}{lcc}
\toprule
\textbf{Member} & \textbf{Primary Responsibilities} & \textbf{Estimated \%} \\
\midrule
{[Student 1]} & [e.g., Literature review, report writing] & \%  \\
{[Student 2]}& [e.g., Model implementation, experiments] & \%  \\
{[Student 3]} & [e.g., Data pipeline, evaluation code, README] & \%  \\
\bottomrule
\end{tabular}

\section{References}

\begin{guidancebox}
\textbf{Guidance.} Use consistent citation formatting. Include at minimum the 5
papers from Section~2, plus any other papers you reference.
arXiv IDs and DOIs are preferred over URLs.
\end{guidancebox}

\begin{enumerate}[label={[\arabic*]},leftmargin=*]
  \item [Author(s), Year. ``Title.'' \textit{Venue}. arXiv:XXXX.XXXXX or DOI.]
  \item [Author(s), Year. ``Title.'' \textit{Venue}.]
  \item [\ldots]
\end{enumerate}

\end{document}
